\section{Entropy of a connected cut}%
\label{sec:peps_connected_cut_entropy}

For a connected region with connected complement, the reduced state of a
regular $G$-isometric network on the region has a flat spectrum, and its rank
is fixed by the number of bonds that cross the cut.

\begin{lemma}[Entropy of the virtual regular boundary]%
    \label{thm:peps_regular_boundary_entropy}
    \lean{TNLean.PEPS.vonNeumannEntropy_regularBoundaryDensity,
        TNLean.PEPS.vonNeumannEntropy_regularBoundaryDensity_eq_mul_log,
        TNLean.PEPS.vonNeumannEntropy_regularBoundaryDensity_eq_boundaryLength}
    \leanok
    Before identifying the virtual
    boundary with a physical PEPS region in
    \cite[Theorem~6.9]{Schuch2010PEPS}, the normalized virtual reduced
    density operator for $b=n+1$ regular boundary legs has von Neumann entropy
    \begin{align}
        S(\rho)=\log(|G|^n)=n\log|G|=b\log|G|-\log|G|.
    \end{align}
\end{lemma}
\begin{proof}\leanok
    \uses{thm:peps_regular_boundary_density, thm:peps_flat_density_entropy}
    The identity $\rho^2=r^{-1}\rho$ applied to each eigenvector gives
    eigenvalue $0$ or $r^{-1}$. With $0\log0=0$, each entropy summand
    is therefore its eigenvalue times $\log r$. Their sum is
    $\operatorname{tr}\rho\log r=\log r$. Since $r=|G|^n$ and $|G|>0$,
    the logarithm laws give the stated expressions.
\end{proof}

\begin{lemma}[Physical realization of a regular boundary state]%
    \label{thm:peps_regular_boundary_physical_realization}
    \lean{TNLean.PEPS.physicalRegularBoundaryState_eq_kronecker_mulVec,
        TNLean.PEPS.physicalRegularBoundarySchmidtMatrix,
        TNLean.PEPS.physicalRegularBoundaryState,
        TNLean.PEPS.physicalRegularBoundaryProjector,
        TNLean.PEPS.physicalRegularBoundaryDensity,
        TNLean.PEPS.physicalRegularBoundarySchmidtMatrix_mul_conjTranspose,
        TNLean.PEPS.partialTrace_physicalRegularBoundaryState,
        TNLean.PEPS.physicalRegularBoundaryProjector_conjTranspose,
        TNLean.PEPS.physicalRegularBoundaryProjector_mul_self,
        TNLean.PEPS.physicalRegularBoundaryProjector_isStarProjection,
        TNLean.PEPS.rank_physicalRegularBoundaryProjector,
        TNLean.PEPS.trace_physicalRegularBoundaryProjector,
        TNLean.PEPS.physicalRegularBoundaryDensity_posSemidef,
        TNLean.PEPS.trace_physicalRegularBoundaryDensity,
        TNLean.PEPS.rank_physicalRegularBoundaryDensity,
        TNLean.PEPS.physicalRegularBoundaryDensity_mul_self,
        TNLean.PEPS.spectrum_physicalRegularBoundaryDensity_subset,
        TNLean.PEPS.vonNeumannEntropy_physicalRegularBoundaryDensity}
    \leanok
    Let $A:\mathbb C[G^b]\to\mathcal H_A$ and
    $B:\mathbb C[G^b]\to\mathcal H_B$ satisfy $A^\dagger A=B^\dagger B=P$.
    Applying $A\otimes B$ to the normalized virtual boundary vector gives
    coefficient matrix $ACB^{\mathsf T}$ and reduced density
    $\rho_A=APA^\dagger/r$. The matrix $Q=APA^\dagger$ is an orthogonal
    projector of rank $r$. Thus $\rho_A$ is positive, has trace one and
    rank $r$, and $\rho_A^2=r^{-1}\rho_A$; its spectrum is contained in
    $\{0,r^{-1}\}$, and $S(\rho_A)=n\log|G|$.

    This proves the preservation of the boundary spectrum under specified
    maps isometric on the invariant subspace. Applying it to a physical
    PEPS region still requires identification of these maps and this
    boundary vector with the geometric contraction in
    \cite[Theorem~6.9]{Schuch2010PEPS}.
\end{lemma}
\begin{proof}\leanok
    \uses{thm:peps_regular_boundary_density, thm:peps_flat_density_entropy}
    The coefficient matrix transforms by $C\mapsto ACB^{\mathsf T}$.
    Since $P$ is real symmetric, $B^\dagger B=P$ gives
    $B^{\mathsf T}\overline B=P$, so the partial trace is
    $ACPC^\dagger A^\dagger=APA^\dagger/r$.
    The equation $A^\dagger A=P$ and $P^2=P$ show $Q^2=Q$ and
    $Q=Q^\dagger$. The map $A$ preserves inner products on
    $\operatorname{ran}P$, and is therefore injective there; its image
    is $\operatorname{ran}Q$, of dimension $r$. The trace and spectrum
    follow from the projection property, and
    Lemma~\ref{thm:peps_flat_density_entropy} gives the entropy.
\end{proof}

\begin{lemma}[Normalization of specified physical cut maps]%
    \label{thm:peps_regular_boundary_normalization}
    \lean{TNLean.PEPS.normalizedRegularBoundaryMap_conjTranspose_mul,
        TNLean.PEPS.normalizedRegularBoundaryMap,
        TNLean.PEPS.physicalRegularBoundarySchmidtMatrix_normalizedRegularBoundaryMap,
        TNLean.PEPS.vonNeumannEntropy_normalizedRegularBoundaryState}
    \leanok
    \uses{thm:peps_regular_boundary_physical_realization,
        thm:peps_regular_open_region_projector}
    Let $M^\dagger M=c_AP$ and $N^\dagger N=c_BP$, with $c_A,c_B>0$.
    Then $\widehat M=M/\sqrt{c_A}$ and
    $\widehat N=N/\sqrt{c_B}$ are isometric on the invariant boundary
    subspace. The normalized bipartite vector with coefficient matrix
    \begin{align}
        \widehat M C\widehat N^{\mathsf T}
        =\frac{MPN^{\mathsf T}}{\sqrt{c_A}\sqrt{c_B}\sqrt r}
        \label{eq:peps_regular_boundary_normalized_matrix}
    \end{align}
    has reduced von Neumann entropy $n\log|G|$.
    For the cut maps of
    Lemma~\ref{thm:peps_regular_open_region_projector}, the numerator
    is the physical PEPS coefficient matrix $K$.

    The Gram identities are explicit assumptions in this statement.
    Lemma~\ref{thm:peps_connected_regular_region_gram} derives them
    from local $G$-isometry for connected untwisted regions, including
    regions containing cycles. The twisted torus contractions of
    \cite[Theorem~6.9]{Schuch2010PEPS} remain separate.
    The scalar convention is recorded in
    \cite{gap:rmp_peps_quantum_double_g_isometry}.
\end{lemma}
\begin{proof}\leanok
    Since $c_A>0$, the product of the normalization scalar and its
    conjugate is $c_A^{-1}$. Thus
    $\widehat M^\dagger\widehat M=P$, and similarly for $N$.
    Pulling the scalars out of the coefficient matrix gives
    (\ref{eq:peps_regular_boundary_normalized_matrix}), and applying
    Lemma~\ref{thm:peps_regular_boundary_physical_realization} to the
    normalized maps yields the entropy.
\end{proof}

\begin{theorem}[Entropy of a connected untwisted finite cut]%
    \label{thm:peps_connected_untwisted_cut_entropy}
    \lean{TNLean.PEPS.regularRegionGramFactor_pos,
        TNLean.PEPS.regularRegionGramFactor,
        TNLean.PEPS.ofReal_regularRegionGramFactor,
        TNLean.PEPS.regularOpenComplementMatrix_eq_regularOpenRegionMatrix,
        TNLean.PEPS.regularPhysicalCutState,
        TNLean.PEPS.exists_positive_gram_regularOpenRegionMatrix_of_connected,
        TNLean.PEPS.exists_normalization_regularPhysicalCutState_of_connected}
    \leanok
    \uses{thm:peps_connected_regular_region_gram,
        thm:peps_regular_boundary_normalization}
    Consider a finite graph PEPS with regular bonds and a $G$-isometric
    tensor at every site. Assume that both $R$ and its complement
    induce nonempty connected graphs, and that $b>0$ bonds cross the cut.
    For the untwisted contraction there is a positive normalization of
    the physical PEPS vector to a unit vector. Its reduced state
    on $R$ has rank $r=|G|^{b-1}$ and a flat spectrum:
    $\rho_R^2=r^{-1}\rho_R$. Its von Neumann entropy is
    \begin{align}
        S(\rho_R)=(b-1)\log|G|.
    \end{align}
    This finite untwisted statement is a restricted consequence of the
    calculation in \cite[Theorem~6.9]{Schuch2010PEPS}; it does not assert
    the theorem for arbitrary torus closures. The remaining scope is
    recorded in \cite{gap:rmp_peps_quantum_double_g_isometry}.
\end{theorem}
\begin{proof}\leanok
    Each local isometry gives a positive Gram factor. Applying
    Lemma~\ref{thm:peps_connected_regular_region_gram} on both sides
    of the cut and normalizing the two maps by their positive factors
    expresses the resulting coefficient matrix as a positive scalar multiple of
    the PEPS cut matrix, so the regular boundary realization gives
    its norm and reduced entropy.
\end{proof}

\begin{lemma}[A local controlled operation removing a boundary twist]%
    \label{thm:peps_regular_boundary_untwisting}
    \lean{TNLean.PEPS.regularTweezerEquiv_apply,
        TNLean.PEPS.regularTweezerEquiv,
        TNLean.PEPS.regularTweezerEquiv_symm_apply,
        TNLean.PEPS.regularTweezer_relative_label,
        TNLean.PEPS.regularTweezerEquiv_symm_apply_twisted,
        TNLean.PEPS.regularTweezerEquiv_symm_apply_conjugated,
        TNLean.PEPS.regularTweezerMatrix,
        TNLean.PEPS.regularTweezerMatrix_mem_unitaryGroup,
        TNLean.PEPS.regularTweezerState,
        TNLean.PEPS.regularTweezerState_factorization,
        TNLean.PEPS.regularTweezerState_mulVec}
    \leanok
    The permutation
    \begin{align}
        F(a,b,c)=(cb^{-1}a,b,c),\qquad
        F^{-1}(a,b,c)=(bc^{-1}a,b,c)
    \end{align}
    acts unitarily in the group basis. If $c=tb$, then
    $F^{-1}(ta,b,tb)=(a,b,tb)$. In particular this removes the
    conjugated twist $t=x^{-1}gx$ from the first label using the
    accessible reference pair $b,c$.

    On the bipartite coefficient vector
    $\mathbf 1_{\{a=td\}}\mathbf 1_{\{c=tb\}}$, applying this inverse
    permutation to the first side gives
    $\mathbf 1_{\{a=d\}}\mathbf 1_{\{c=tb\}}$. All dependence on $t$
    is then confined to the reference pair on that side. This proves
    the local coordinate operation in
    \cite{Schuch2010PEPS}, displayed at local source lines 1961--1965.
    Identifying these coordinates with a whole torus contraction is
    a further step.
\end{lemma}
\begin{proof}\leanok
    Group cancellation verifies the two inverse permutations and
    $cb^{-1}=t$. A permutation of an orthonormal basis is unitary.
    Substitution in the two indicator functions yields the factorization.
\end{proof}
